\problem{}

ShanghaiTech University wants to deploy a night security patrol system on campus. 
Model the campus road network as an undirected graph $G=(V,E)$, where:
\begin{itemize}
    \item each vertex $v \in V$ represents an important campus location (e.g., a building entrance, a road intersection, or a gate);
    \item each edge $e=(u,v)\in E$ represents a road segment connecting two locations.
\end{itemize}
A security guard stationed at location $v$ can patrol (and thus \emph{cover}) all road segments incident to $v$. 
A set $S \subseteq V$ is called a \emph{patrol placement} if every road segment on campus is covered by at least one guard; 
equivalently, for every edge $(u,v)\in E$, at least one of $u$ or $v$ belongs to $S$.

\medskip
Consider the following simple deterministic strategy:
\begin{quote}
\textbf{Start} with $S \leftarrow \emptyset$. \\
\textbf{While} there exists an uncovered road segment $e=(u,v)$:
\begin{itemize}
    \item Add both endpoints to $S$, i.e., $S \leftarrow S \cup \{u,v\}$.
    \item Delete all road segments incident on $u$ or $v$ (since they are now covered).
\end{itemize}
\end{quote}

\noindent \textbf{Task.} 
Show that the above algorithm is a $2$-approximation for the minimum patrol placement problem.
In particular, prove that:
\begin{enumerate}
    \item $S$ is a valid patrol placement (i.e., it covers every road segment in $E$), and
    \item $|S| \le 2|S^*|$, where $S^*$ is an optimal (minimum-size) patrol placement.
\end{enumerate}

\noindent \textbf{Solution:}



Define indicator random variables for each index \(i\):
\[
X_i = \begin{cases}
1, & \text{if position } i \text{ is a local maximum},\\
0, & \text{otherwise}.
\end{cases}
\]
Then \(M = \sum_{i=1}^n X_i\), and by linearity of expectation,
\[
\mathbb{E}[M] = \sum_{i=1}^n \mathbb{E}[X_i] = \sum_{i=1}^n \Pr(X_i = 1).
\]

We compute the probability that position \(i\) is a local maximum.

\noindent\textbf{Boundary positions:}
\begin{itemize}
    \item For \(i = 1\): We require \(A[1] > A[2]\). Since the permutation is uniformly random, \(A[1]\) and \(A[2]\) are two distinct random numbers, and all orderings are equally likely. Hence,
    \[
    \Pr(X_1 = 1) = \frac{1}{2}.
    \]
    \item For \(i = n\): Similarly, \(\Pr(X_n = 1) = \Pr(A[n] > A[n-1]) = \frac{1}{2}\).
\end{itemize}

\noindent\textbf{Internal positions \(2 \leq i \leq n-1\):}
We require \(A[i-1] < A[i]\) and \(A[i] > A[i+1]\). Consider the three consecutive values \(A[i-1], A[i], A[i+1]\). They are three distinct numbers chosen uniformly from \(\{1,\dots,n\}\), and all \(3! = 6\) orderings are equally likely. For \(A[i]\) to be the largest among the three, it must be the maximum. The probability that the middle element is the largest is \(\frac{2}{6} = \frac{1}{3}\). Thus,
\[
\Pr(X_i = 1) = \frac{1}{3}.
\]

Therefore, for \(n \ge 2\),
\[
\mathbb{E}[M] = \frac{1}{2} + \frac{1}{2} + (n-2)\cdot\frac{1}{3} = 1 + \frac{n-2}{3} = \frac{n+1}{3}.
\]

(For \(n = 1\), the single element is trivially a local maximum, so \(\mathbb{E}[M] = 1\), but the formula \(\frac{n+1}{3}\) does not apply.)

Hence, the expected number of local maxima is \({\frac{n+1}{3}}\) for \(n \ge 2\).


\newpage